\label{ch:ew}

\section{Funciones constitutiva y determinantal del operador angular}

Recordemos el operador de dos hojas del \cref{ch:vacuum},
\begin{equation}
\label{eq:ew-angular-operator}
T=q_+P_+ +q_-P_-,
\qquad
q_-=\ee^{-(A_{\rm rad}-C^*_{\rm rad})},
\qquad
q_+=\ee^{-(A_{\rm rad}+C^*_{\rm rad})}.
\end{equation}
El vac\'io electromagn\'etico aplica la funci\'on
\begin{equation}
\label{eq:ew-vacuum-reader}
F(T)^2,
\qquad
F(z)=\frac{1-z^{90}}{1-z^{120}},
\end{equation}
mientras el sector angular electrod\'ebil aplica el determinante
\begin{equation}
\label{eq:ew-determinant-reader}
\boxed{
\det T=q_+q_-=\ee^{-2A_{\rm rad}}=:D_A.}
\end{equation}
Son dos realizaciones funcionales del antecedente com\'un \((A,C^*,T)\). No existe una
flecha causal \(F(T)^2\to\det T\) ni su inversa. La unidad estructural
reside en \(T\), no en identificar despu\'es permeabilidad, permitividad y
\(\theta_W\).

En la involuci\'on de hojas \(C^*\mapsto-C^*\), los autovalores
\(q_+,q_-\) se intercambian. Por ello \(D_A\) y la velocidad constitutiva
son pares, mientras \((\widehat\mu,\widehat\varepsilon)\) se intercambia y
\(\widehat Z\) se invierte. La relaci\'on electrod\'ebil que sigue usa el
invariante par.

\section{Relaci\'on exacta entre los sectores \(W\) y \(Z\)}

La coordenada angular de duraci\'on--perduraci\'on es
\begin{equation}
\label{eq:ew-angular-character}
\chi_{\rm dur}(n_A,s_C)
=\exp\!\left(n_AA_{\rm rad}
+\frac{s_C}{6}C^*_{\rm rad}\right).
\end{equation}
Esta coordenada no se introduce como una etiqueta aislada. Su procedencia es
\begin{equation}
\label{eq:ew-signature-provenance}
\begin{aligned}
\APP&\longrightarrow\TRIT\longrightarrow\TPK
\longrightarrow\text{extensión superviviente}
\longrightarrow\text{ruta}\\
&\longrightarrow\text{ledger}
\longrightarrow(n_A,s_C,\ldots)
\longrightarrow\chi_{\rm dur}.
\end{aligned}
\end{equation}
Las firmas de \(W\), \(Z\) y \(H\) son magnitudes tipadas obtenidas mediante
esa cadena; el presente cap\'itulo las adopta sin seleccionarlas a partir de
sus masas terminales. En la realizaci\'on angular establecida, los sectores
cargado y neutro satisfacen
\begin{equation}
\label{eq:ew-WZ-signatures}
(n_A,s_C)_W=(94,-1),
\qquad
(n_A,s_C)_Z=(95,-1).
\end{equation}
Esta afirmaci\'on se refiere a la componente angular com\'un. Las
coordenadas adicionales de refinamiento de la ley general de masas no se
eliminan salvo que el transductor f\'isico demuestre su cancelaci\'on.

\begin{theorem}[Defecto de Weinberg en el esquema de masas f\'isicas de la carta angular]
\label{thm:ew-weinberg}
Si \(W\) y \(Z\) se eval\'uan en una misma secci\'on de escala y comparten
el refinamiento que acompa\~na a \cref{eq:ew-WZ-signatures}, entonces
\begin{equation}
\label{eq:ew-WZ-ratio}
\boxed{
\frac{m_W^{\angle}}{m_Z^{\angle}}
=\ee^{-A_{\rm rad}}=\sqrt{D_A},
\qquad
\left(\frac{m_W^{\angle}}{m_Z^{\angle}}\right)^2=D_A.}
\end{equation}
En el esquema de masas f\'isicas,
\begin{equation}
\label{eq:ew-sin2thetaW}
\boxed{
\sin^2\theta_W^{\rm OS,HMT}=1-D_A
=0.2248708807528435698111159035\ldots .}
\end{equation}
Adem\'as,
\begin{equation}
\label{eq:ew-WZ-terminal-ratio}
\frac{m_W^{\angle}}{m_Z^{\angle}}
=0.8804141748331613670128885459\ldots .
\end{equation}
\end{theorem}

\begin{proof}
Por \cref{eq:ew-angular-character,eq:ew-WZ-signatures},
\[
\frac{\chi_{\rm dur}(94,-1)}{\chi_{\rm dur}(95,-1)}
=\exp(-A_{\rm rad});
\]
la coordenada \(C^*\) se cancela porque ambos sectores comparten
\(s_C=-1\). Al elevar al cuadrado y usar
\cref{eq:ew-determinant-reader} se obtiene \(D_A\). La definici\'on
de masas f\'isicas \(s_W^2=1-m_W^2/m_Z^2\) produce
\cref{eq:ew-sin2thetaW}.
\end{proof}

El contenido principal es la identidad entre dos construcciones
internos: el determinante del operador angular y la diferencia de un paso
en la realizaci\'on \(W/Z\). La promoci\'on a selector f\'isico de calibre exige
un entrelazador que lleve las hojas de \(T\) a la base neutra
\((W^3,B)\), y que demuestre qu\'e refinamientos se conservan o cancelan.

\section{Acoplamientos de calibre en normalizaci\'on racionalizada}

Definimos el acoplamiento adimensional de calibre
\begin{equation}
\label{eq:ew-egauge}
e_g:=\sqrt{4\pi\alpha_{\HMT}}
=0.3028221208717526427662442689\ldots .
\end{equation}
No debe confundirse \(e_g\) con una secci\'on dimensional de carga
\(e_{\rm int}\): la primera es un acoplamiento en unidades naturales; la
segunda vive en la recta de carga de la Mec\'anica Dimensional.

Sea
\begin{equation}
\label{eq:ew-sw-cw}
s_W^2=1-D_A,
\qquad
c_W^2=D_A.
\end{equation}
La convenci\'on electrod\'ebil racionalizada a nivel de \`arbol da
\begin{align}
g&=\frac{e_g}{s_W}
=2\sqrt{\frac{\pi\alpha_{\HMT}}{1-D_A}}
=0.6385883427334574283647003809\ldots,
\label{eq:ew-g}\\
g'&=\frac{e_g}{c_W}
=2\sqrt{\frac{\pi\alpha_{\HMT}}{D_A}}
=0.3439541633108502429362319257\ldots,
\label{eq:ew-gprime}\\
\frac{g'}{g}
&=\sqrt{D_A^{-1}-1}
=0.5386164141966093594996610933\ldots .
\label{eq:ew-gprime-over-g}
\end{align}

\begin{proposition}[Relaci\'on custodial de \`arbol]
\label{prop:ew-rho}
En la misma interfaz,
\begin{equation}
\label{eq:ew-rho}
\rho_{\rm EW}^{(0)}
:=\frac{(m_W^{\angle})^2}
{(m_Z^{\angle})^2c_W^2}=1.
\end{equation}
\end{proposition}

\begin{proof}
Por \cref{eq:ew-WZ-ratio,eq:ew-sw-cw}, tanto el numerador relativo como
\(c_W^2\) son \(D_A\).
\end{proof}

Las ecuaciones \cref{eq:ew-g,eq:ew-gprime} son composiciones exactas una
vez declarados el esquema de masas f\'isicas, la racionalizaci\'on y la escala. No
afirman que un cambio de esquema de renormalizaci\'on deje inalterados sus
valores num\'ericos.

\section{Constante de Fermi en aproximaci\'on de \`arbol}

En unidades naturales, la integraci\'on del propagador cargado a baja
energ\'ia produce
\begin{equation}
\label{eq:ew-GF-def}
\frac{G_F^{(0)}}{\sqrt2}=\frac{g^2}{8m_W^2}.
\end{equation}
Sustituyendo \cref{eq:ew-g} se obtiene la composici\'on
\begin{equation}
\label{eq:ew-GF}
\boxed{
G_F^{(0)}
=\frac{\pi\alpha_{\HMT}}
{\sqrt2(1-D_A)m_W^2}.}
\end{equation}
Al adoptar la masa obtenida previamente
\begin{equation}
\label{eq:ew-W-output}
m_W^{\HMT}=80.381592550221\ldots\ \mathrm{GeV},
\end{equation}
la evaluaci\'on terminal es
\begin{equation}
\label{eq:ew-GF-value}
\boxed{
G_F^{(0)}
=1.1157162817659203186621784\times10^{-5}\ \mathrm{GeV}^{-2}.}
\end{equation}

La diferencia respecto de una constante de Fermi renormalizada no se usa
para elegir ning\'un coeficiente. En la convenci\'on
\begin{equation}
\label{eq:ew-Delta-r}
G_F
=\frac{\pi\alpha_{\HMT}}
{\sqrt2(1-D_A)m_W^2(1-\Delta r)},
\end{equation}
la obligaci\'on restante es derivar \(\Delta r\) desde el sector interno de
bucles, estados y escala. Abrir el valor observado de \(G_F\) para definir
\(\Delta r\) ser\'ia una reconstrucci\'on calibrada, no una predicci\'on.

\section{Dependencia orientada del acoplamiento de Higgs}

En la carta angular, las rutas pertinentes son
\begin{equation}
\label{eq:ew-HW-signatures}
(n_A,s_C)_H=(97,9),
\qquad
(n_A,s_C)_W=(94,-1).
\end{equation}
Por tanto,
\begin{equation}
\label{eq:ew-HW-ratio}
\boxed{
\frac{m_H^{\angle}}{m_W^{\angle}}
=\frac{\chi_{\rm dur}(97,9)}{\chi_{\rm dur}(94,-1)}
=\exp\!\left(3A_{\rm rad}+\frac53C^*_{\rm rad}\right)
=1.55872087212325548564\ldots .}
\end{equation}
Aqu\'i \(C^*\) no se cancela: el sector escalar conserva la orientaci\'on
que el determinante \(D_A=\ee^{-2A_{\rm rad}}\) no distingue.

Tomemos la convenci\'on del potencial
\begin{equation}
\label{eq:ew-Higgs-potential}
V(H)=-\mu_H^2H^\dagger H
+\lambda_H(H^\dagger H)^2,
\end{equation}
para la cual
\begin{equation}
\label{eq:ew-Higgs-tree-relations}
m_H^2=2\lambda_Hv^2,
\qquad
m_W^2=\frac{g^2v^2}{4}.
\end{equation}

\begin{corollary}[Autoacoplamiento orientado de Higgs]
\label{cor:ew-Higgs-lambda}
En la interfaz de \`arbol declarada,
\begin{equation}
\label{eq:ew-Higgs-lambda}
\boxed{
\lambda_H
=\frac{\pi\alpha_{\HMT}}{2(1-D_A)}
\exp\!\left(6A_{\rm rad}+\frac{10}{3}C^*_{\rm rad}\right)
=0.1238479115482466804662\ldots .}
\end{equation}
Adem\'as,
\begin{equation}
\label{eq:ew-GF-Higgs}
G_F^{(0)}m_H^2=\sqrt2\,\lambda_H.
\end{equation}
\end{corollary}

\begin{proof}
Dividir las dos ecuaciones de
\cref{eq:ew-Higgs-tree-relations} da
\(m_H^2/m_W^2=8\lambda_H/g^2\). Con
\(g^2=4\pi\alpha_{\HMT}/(1-D_A)\) y el cuadrado de
\cref{eq:ew-HW-ratio} se obtiene \cref{eq:ew-Higgs-lambda}. Por otra parte,
\(G_F^{(0)}=1/(\sqrt2v^2)\) y
\(m_H^2=2\lambda_Hv^2\) implican
\cref{eq:ew-GF-Higgs}.
\end{proof}

El valor de \(\lambda_H\) es anterior a la incorporaci\'on de correcciones
radiativas y depende de
la convenci\'on y la escala declaradas. Su inter\'es estructural es que
separa dos memorias: el sector com\'un de calibre emplea el invariante par \(D_A\),
mientras la ruta de Higgs retiene expl\'icitamente \(C^*\).

\section{Clausura determinantal del sector electrodébil: mezcla, acoplamientos y espectro constitutivo}

La cadena completa de este cap\'itulo es
\begin{equation}
\label{eq:ew-causal-diagram}
\boxed{
\begin{aligned}
(A,C^*)&\longrightarrow T
\longrightarrow
\begin{cases}
F(T)^2&\longrightarrow
(\widehat\mu,\widehat\varepsilon,\widehat Z,\widehat c),\\
\det T=D_A&\longrightarrow
(s_W^2,c_W^2),
\end{cases}\\
(D_A,\alpha_{\HMT})&\longrightarrow(e_g,g,g'),\\
(D_A,\alpha_{\HMT},m_W)&\longrightarrow G_F^{(0)},\\
(A,C^*,D_A,\alpha_{\HMT})&\longrightarrow
\left(\frac{m_H}{m_W},\lambda_H\right).
\end{aligned}}
\end{equation}
La cifra aparece siempre al final. Ni \(\sin^2\theta_W\), ni
\(G_F\), ni \(\lambda_H\) se introducen como valores objetivo para seleccionar
\(A,C^*\) o las firmas de ruta.

\begin{proposition}[Resultado principal]
El mismo operador angular determina el vac\'io mediante una funci\'on racional
de sus hojas y el defecto electrod\'ebil mediante su determinante. La
diferencia de un paso entre \(W\) y \(Z\) convierte ese determinante en
\(\sin^2\theta_W^{\rm OS}=1-D_A\); la ruta de Higgs, en cambio, conserva
el contra\'angulo y conserva la orientaci\'on eliminada por el determinante.
\end{proposition}

\begin{remark}[Control negativo]
La relaci\'on de Weinberg queda refutada en su dominio si las firmas
angulares \(W/Z\) no difieren s\'olo por \(n_A:94\to95\), si una coordenada
de refinamiento no se cancela en la interfaz declarada o si
\((m_W^{\angle}/m_Z^{\angle})^2\neq D_A\). La promoci\'on de calibre queda
bloqueada si no existe un entrelazador con la base \((W^3,B)\). Las cifras
de \(g,g',G_F,\lambda_H\) cambian si se altera esquema o escala sin
registrarlo. Usar \(G_F\) o \(m_H\) observados para seleccionar
\(\Delta r\), las firmas o \(C^*\) introduce retroalimentaci\'on por el
valor objetivo e invalida la derivaci\'on. Confundir \(e_g\) con la carga
dimensional constituye un error de tipo.
\end{remark}

\subsection*{Alcance lógico y dominio de validez}
Determinante angular, cociente \(W/Z\) y defecto complementario de Weinberg:
\emph{Teorema interno exacto en la carta angular condicionada}. Acoplamientos \(g,g'\),
\(\rho_{\rm EW}^{(0)}\), Fermi y Higgs: \emph{Composici\'on exacta en la
interfaz electrod\'ebil de \`arbol}. El entrelazador de base de calibre y la
correcci\'on interna \(\Delta r\): \emph{Frontera abierta con criterio de refutación}.
Fermi y Higgs no son constantes independientes seleccionadas sin valores
objetivo: son composiciones exactas que emplean las firmas, la magnitud
\(m_W\) y la
convenci\'on de \`arbol declarada. Los valores convencionales posteriores
son \emph{Contraste metrol\'ogico externo}.

